\exercisesection{2}

\begin{enumerate}
\def\labelenumi{\arabic{enumi}.}
\item
  给出左 A-模的正合序列 \(0 \rightarrow N' \rightarrow N \rightarrow N'' \rightarrow 0\) 与右 A-模 E 的一个例子，使 E 是 \(N'\) -平坦与 \(N''\) -平坦的，但不是 N-平坦的（例如取 \(N' = N'' = Z/2Z\) ）。
\item
  设 M、N 是 A-模 E 的两个子模，且 \(M + N\) 是平坦的。为使 M 与 N 都平坦，必须且只需 \(M \cap N\) 是平坦的。
\item
  设 A 是域 K 上两个未定元的多项式环 K{[}X, Y{]} 。
\end{enumerate}

\begin{enumerate}
\def\labelenumi{(\alph{enumi})}
\item
  在 A 中考虑主理想 \(\mathfrak{b} = (\mathbf{X}), \mathfrak{c} = (\mathbf{Y})\) ，它们是自由 A-模，且其交 \(\mathfrak{b} \cap \mathfrak{c} = (\mathbf{XY})\) 也是自由的。证明 \(\mathfrak{a} = \mathfrak{b} + \mathfrak{c}\) 不是平坦 A-模，尽管 \(\mathfrak{a}\) 是无挠的（参看《代数学》第三章 §2 习题 4）。
\item
  在 A-模 \(A^{2}\) 中，设 R 是由元素 \((x, -x)\) （其中 \(x \in a\) ）组成的子模。在 A-模 \(A^{2}/R\) 中，设 M、N 是 \(A^{2}\) 的因子子模的像所成的子模；证明 M 与 N 都同构于 A，但 \(M \cap N\) 不是平坦 A-模。
\end{enumerate}

\begin{enumerate}
\def\labelenumi{\arabic{enumi}.}
\setcounter{enumi}{3}
\tightlist
\item
  \begin{enumerate}
  \def\labelenumii{(\alph{enumii})}
  \tightlist
  \item
    给出正合序列的一个例子
  \end{enumerate}
\end{enumerate}

\[
0 \rightarrow \mathrm{E} ^ {\prime} \rightarrow \mathrm{E} \rightarrow \mathrm{E} ^ {\prime \prime} \rightarrow 0
\]

它不分裂，且它的所有项都是平坦模（参看《代数学》第七章 §3 习题 8 (b)）。

\begin{enumerate}
\def\labelenumi{(\alph{enumi})}
\setcounter{enumi}{1}
\tightlist
\item
  由 (a) 导出正合序列的一个例子
\end{enumerate}

\[
0 \rightarrow \mathrm{E} ^ {\prime} \rightarrow \mathrm{E} \rightarrow \mathrm{E} ^ {\prime \prime} \rightarrow 0
\]

它不分裂，且它的各项都是非平坦的右 A-模，使得对每个左 A-模 F，序列

\[
0 \rightarrow \mathrm{E} ^ {\prime} \otimes \mathrm{F} \rightarrow \mathrm{E} \otimes \mathrm{F} \rightarrow \mathrm{E} ^ {\prime \prime} \otimes \mathrm{F} \rightarrow 0
\]

是正合的（利用第 1 目的引理 2）。

\begin{enumerate}
\def\labelenumi{\arabic{enumi}.}
\setcounter{enumi}{4}
\tightlist
\item
  给出右 A-模 E、左 A-模 F 以及 F 的两个子模 \(F'\) 、 \(F''\) 的一个例子，使 \(E \otimes (F' \cap F'')\) 在 \(E \otimes F\) 中的典范像不是 \(E \otimes F'\) 与 \(E \otimes F''\) 的典范像的交（参看习题 3(a)）。
\end{enumerate}

¶ 6. 设 A 是环，M 是左 A-模。正合序列

\[
\mathrm{L} _ {n} \rightarrow \mathrm{L} _ {n - 1} \rightarrow \dots \rightarrow \mathrm{L} _ {1} \rightarrow \mathrm{L} _ {0} \rightarrow \mathrm{M} \rightarrow 0
\]

其中 \(L_{i}\) 是自由左 A-模（ \(0 \leqslant i \leqslant n\) ），称为 M 的长度为 n 的表示或 M 的 n-表示。若所有 \(L_{i}\) 都是有限生成的自由模，则称该表示为有限表示。

若 M 是有限生成的左 A-模，用 \(\lambda(M)\) 表示一个上确界（有限的或 \(+\infty\) ），即整数 \(n \geqslant 0\) （M 对之有有限 \(n\) -表示）的上确界。若 M 不是有限生成的，则置 \(\lambda(M) = -1\) 。

\begin{enumerate}
\def\labelenumi{(\alph{enumi})}
\item
  设 \(0 \to \mathbf{P} \to \mathbf{N} \to \mathbf{M} \to 0\) 是左 A-模的正合序列。则 \(\lambda(\mathbf{N}) \geqslant \inf(\lambda(\mathbf{P}), \lambda(\mathbf{M}))\) 。（从两个 \(n\) -表示——分别属于 \(\mathbf{P}\) 与 \(\mathbf{M}\) ——出发，利用 §1 习题 5 为 \(\mathbf{N}\) 导出一个表示。）
\item
  设 \(M_{n} \xrightarrow{u_{n}} M_{n-1} \rightarrow \cdots \rightarrow M_{0} \xrightarrow{u_{0}} M \rightarrow 0\) 是 M 的有限 n-表示；证明：若 \(\lambda(M) > n\) ，则 \(\operatorname{Ker}(u_{n})\) 是有限生成的 A-模。（设
\end{enumerate}

\[
\mathrm{L} _ {n + 1} \xrightarrow {v _ {n + 1}} \mathrm{L} _ {n} \longrightarrow \dots \longrightarrow \mathrm{L} _ {2} \xrightarrow {v _ {2}} \mathrm{L} _ {1} \xrightarrow {v _ {1}} \mathrm{L} _ {0} \xrightarrow {v _ {0}} \mathrm{M} \longrightarrow 0
\]

是一个有限 \((n + 1)\) -表示——它属于 \(\mathbf{M}\) ；令 \(\mathbf{P} = \operatorname{Ker}(v_0)\) ，使存在一个 \(n\) -表示，属于 \(\mathbf{P}\) ：

\[
\mathrm{L} _ {n + 1} \xrightarrow {v _ {n + 1}} \mathrm{L} _ {n} \longrightarrow \dots \longrightarrow \mathrm{L} _ {2} \xrightarrow {v _ {2}} \mathrm{L} _ {1} \xrightarrow {v _ {1}} \mathrm{P} \longrightarrow 0.
\]

把 (a) 的方法应用于正合序列 \(0 \rightarrow P \rightarrow L_{0} \rightarrow M \rightarrow 0\) ，得到一个正合序列

\[
\mathrm{M} _ {n} \oplus \mathrm{L} _ {n + 1} \xrightarrow {w _ {n}} \mathrm{M} _ {n - 1} \oplus \mathrm{L} _ {n} \xrightarrow {w _ {n - 1}} \dots \xrightarrow {w _ {1}} \mathrm{M} _ {0} \oplus \mathrm{L} _ {1} \xrightarrow {w _ {0}} \mathrm{L} _ {0} \longrightarrow 0
\]

以及正合序列 \(0 \to \operatorname{Ker}(v_{i+1}) \to \operatorname{Ker}(w_i) \to \operatorname{Ker}(u_i) \to 0\) （§1 第 4 目，命题 2）。最后注意 \(\operatorname{Ker}(w_i)\) 是 \(\mathbf{M}_i \oplus \mathbf{L}_{i+1}\) 的直因子。

\begin{enumerate}
\def\labelenumi{(\alph{enumi})}
\setcounter{enumi}{2}
\tightlist
\item
  证明在 (a) 的假设下
\end{enumerate}

\[
\lambda (\mathbf {M}) \geqslant \inf (\lambda (\mathbf {N}), \lambda (\mathbf {P}) + 1).
\]

（若 \(n \leqslant \inf(\lambda(N), \lambda(P) + 1)\) ，则对 \(n\) 用归纳法并如 (a) 中那样论证且利用 (b)，证明 \(\lambda(M) \geqslant n\) 。）

\begin{enumerate}
\def\labelenumi{(\alph{enumi})}
\setcounter{enumi}{3}
\tightlist
\item
  证明在 (a) 的假设下
\end{enumerate}

\[
\lambda (\mathrm{P}) \geqslant \inf (\lambda (\mathrm{N}), \lambda (\mathrm{M}) - 1).
\]

（与 (c) 中方法相同。）由此推出：若 \(\lambda(N) = +\infty\) ，则 \(\lambda(M) = \lambda(P) + 1\) 。

\begin{enumerate}
\def\labelenumi{(\alph{enumi})}
\setcounter{enumi}{4}
\tightlist
\item
  由 (a)、(c) 与 (d) 推出：若 \(\mathbf{N} = \mathbf{M} \oplus \mathbf{P}\) ，则
\end{enumerate}

\[
\lambda (\mathbf {N}) = \inf (\lambda (\mathbf {M}), \lambda (\mathbf {P})).
\]

特别地，为使 N 有有限表示，必须且只需 M 与 P 也都有有限表示。

\begin{enumerate}
\def\labelenumi{(\alph{enumi})}
\setcounter{enumi}{5}
\tightlist
\item
  设 \(N_{1}\) 、 \(N_{2}\) 是 A-模 M 的两个子模。假设 \(N_{1}\) 与 \(N_{2}\) 都有有限表示。为使 \(N_{1} + N_{2}\) 有有限表示，必须且只需 \(N_{1} \cap N_{2}\) 是有限生成的。
\end{enumerate}

\begin{enumerate}
\def\labelenumi{\arabic{enumi}.}
\setcounter{enumi}{6}
\tightlist
\item
  \begin{enumerate}
  \def\labelenumii{(\alph{enumii})}
  \tightlist
  \item
    用习题 6 的记号，证明：若 M 是投射模，则 \(\lambda(M) = -1\) 或 \(\lambda(M) = +\infty\) 。若 A 是左 Noether 环，则对每个 A-模 M， \(\lambda(M) = -1\) 或 \(\lambda(M) = +\infty\) 。
  \end{enumerate}
\end{enumerate}

\begin{enumerate}
\def\labelenumi{(\alph{enumi})}
\setcounter{enumi}{1}
\item
若 \(\mathfrak{a}\) 是环 \(\mathbf{A}\) 的一个非有限生成的左理想，则 \(\mathbf{A}_s / \mathfrak{a}\) 是单生成 A-模，且不具有有限表示，换言之 \(\lambda(\mathbf{A}_s / \mathfrak{a}) = 0\) （第 8 目，引理 9）。
\item
给出环 A 的单生成左理想 \(\mathfrak{a}\) 的一个例子，使 \(\mathbf{A}_s / \mathfrak{a}\) （它有有限表示）的对偶不是有限生成的右 A-模。
\item
设 K 是交换域，E 是向量空间 \(\mathbf{K}^{(\mathbf{N})}, (e_n)\) （第二个记号是它的典范基），T 是 E 的张量代数，因此它有一个由有限乘积 \(e_{i_1}e_{i_2}\ldots e_{i_k}\) ( \(k \geqslant 0, i_j \in N\) ，对所有 j) 组成的基。对给定的整数 n，设 b 是 T 中由乘积
\end{enumerate}

\[
e _ {1} e _ {0}, e _ {2} e _ {1}, \dots , e _ {n} e _ {n - 1}
\]

与 \(e_{n+k}e_n\) （对所有 \(k\geqslant 1\) ）生成的双边理想；设 A 是商环 T/b，且对每个整数 \(m\) ，设 \(a_m\) 是 \(e_m\) 在 A 中的典范像。证明：若 \(\mathbf{M} = \mathbf{A}_s / \mathbf{A}a_0\) ，则 \(\lambda (\mathbf{M}) = n\) （注意：对 \(m\leqslant n - 1\) ， \(a_m\) 的左零化子是 \(\mathrm{A}a_{m+1}\) ，并利用习题 6(b)）。

\begin{enumerate}
\def\labelenumi{\arabic{enumi}.}
\setcounter{enumi}{7}
\item
设 C 是交换环，E、F 是两个 C-模。证明 \(\lambda(E \otimes_{\mathbb{C}} F) \geqslant \inf(\lambda(E), \lambda(F))\) 。
\item
设 E 是有限表示的左 A-模。
\end{enumerate}

\begin{enumerate}
\def\labelenumi{(\alph{enumi})}
\item
证明对右 A-模的每个族 \((\mathbf{F}_t)_{t \in I}\) ，典范同态 \(\mathbf{E} \otimes_{\mathbf{A}} \left( \prod_{t \in I} \mathbf{F}_t \right) \to \prod_{t \in I} (\mathbf{E} \otimes_{\mathbf{A}} \mathbf{F}_t)\) （《代数学》第二章 §3 第 7 目）是双射。
\item
设 \((G_{\alpha}, \phi_{\beta\alpha})\) 是左 A-模的一个归纳系；证明典范同态
\end{enumerate}

\[
\varinjlim \operatorname{Hom} _ {\mathrm{A}} (\mathrm{E}, \mathrm{G} _ {\alpha}) \rightarrow \operatorname{Hom} _ {\mathrm{A}} (\mathrm{E}, \varinjlim \mathrm{G} _ {\alpha})
\]

是双射。

\begin{enumerate}
\def\labelenumi{\arabic{enumi}.}
\setcounter{enumi}{9}
\tightlist
\item
  \begin{enumerate}
  \def\labelenumii{(\alph{enumii})}
  \tightlist
  \item
  设 A 是环，I 是集合，R 是 \(L = A_{d}^{(I)}\) 的子模。设 S 是有序偶 (J, S) 的集合，其中 J 是 I 的有限子集，S 是 \(A_{d}^{J} \cap R\) 的有限生成子模。用关系``J ⊂ J' 且 S ⊂ S'\,''给 S 排序；证明 S 关于此序关系是定向的，族 ( \(A_{d}^{J}/S\) ) 是以 S 为指标集的右 A-模的归纳系，且存在 L/R 到 \(\varinjlim (A_{d}^{J}/S)\) 上的同构。
  \end{enumerate}
\end{enumerate}

\begin{enumerate}
\def\labelenumi{(\alph{enumi})}
\setcounter{enumi}{1}
\tightlist
\item
由 (a) 推出：每个 A-模都是有限表示 A-模的归纳极限。
\end{enumerate}

\begin{enumerate}
\def\labelenumi{\arabic{enumi}.}
\setcounter{enumi}{10}
\tightlist
\item
设 E 是右 A-模。若 E 的每个有限生成子模都是有限表示的，则称 E 为伪凝聚模；伪凝聚模的每个子模都是伪凝聚的。若 E 是伪凝聚的且有限生成的（从而是有限表示的），则称 E 为凝聚模。
\end{enumerate}

\begin{enumerate}
\def\labelenumi{(\alph{enumi})}
\item
设 \(0 \to E' \to E \to E'' \to 0\) 是右 A-模的正合序列。证明：若 E 是伪凝聚（相应地凝聚）的且 \(E'\) 有限生成，则 \(E''\) 是伪凝聚（相应地凝聚）的。证明：若 \(E'\) 与 \(E''\) 都伪凝聚（相应地凝聚），则 E 亦然。证明：若 E 与 \(E''\) 都凝聚，则 \(E'\) 亦然（利用习题 6 与第 8 目的引理 9）。
\item
设 E 是凝聚 A-模， \(E'\) 是伪凝聚（相应地凝聚）A-模。证明：对每个同态 \(u: E \to E'\) ， \(\text{Im}(u)\) 与 \(\text{Ker}(u)\) 是凝聚的，且 \(\text{Coker}(u)\) 是伪凝聚（相应地凝聚）的（利用 (a)）。
\item
证明伪凝聚（相应地凝聚）模的每个直和（相应地每个有限直和）是伪凝聚（相应地凝聚）模。
\item
若 E 是伪凝聚模，M、N 是 E 的凝聚子模，证明 \(M + N\) 与 \(M \cap N\) 是凝聚的（利用 (a) 与 (c)）。
\item
假设 A 交换。证明：若 E 是凝聚 A-模且 F 是凝聚（相应地伪凝聚）A-模，则 \(\operatorname{Hom}_{\mathbf{A}}(\mathbf{E}, \mathbf{F})\) 是凝聚（相应地伪凝聚）A-模。（归结为 F 为凝聚的情形，考虑 E 的一个有限表示，然后利用 (b)。）
\end{enumerate}

¶ 12. (a) 设 A 是环。证明下列四个性质等价：

\((\alpha)\) 右 A-模 \(\mathbf{A}_d\) 是凝聚的（习题 11）。

(β) 每个有限表示的右 A-模都是凝聚的。

(γ) 每个 A-模 \(A_{s}^{I}\) （I 为任意集合）都是平坦的。

(δ) 平坦左 A-模的每个乘积都是平坦的。

（为证明 \((\alpha)\) 蕴含 \((\beta)\) ，利用习题 11(b)。为看出 \((\gamma)\) 蕴含 \((\alpha)\) ，利用第 11 目的命题 13 并用反证法论证。为证明 \((\alpha)\) 蕴含 \((\delta)\) ，利用习题 9。）

这样的环称为右凝聚环，左凝聚环的概念类似地定义。

\begin{enumerate}
\def\labelenumi{(\alph{enumi})}
\setcounter{enumi}{1}
\tightlist
\item
  证明每个右 Noether 环都是右凝聚的。给出一个不是左凝聚的右 Artin 环的例子（参看《代数学》第八章 §2 习题 4）。
\end{enumerate}

*(c) 秩 1 的非离散赋值环（第六章）是凝聚的，但它含有非凝聚的理想，并且有（单生成的）非伪凝聚的商模。*

\begin{enumerate}
\def\labelenumi{(\alph{enumi})}
\setcounter{enumi}{3}
\item
  证明：若 A 是右凝聚环，则对每个右 A-模 E， \(\lambda(\mathrm{E}) = -1\) 或 \(\lambda(\mathrm{E}) = 0\) 或 \(\lambda(\mathrm{E}) = +\infty\) （习题 6）。
\item
设 \((\mathbf{A}_{\alpha}, \phi_{\beta \alpha})\) 是环的归纳系，其指标集是定向的，且设 \(\mathbf{A} = \varinjlim \mathbf{A}_{\alpha}\) 。假设对 \(\alpha \leqslant \beta\) ， \(\mathbf{A}_{\beta}\) 是平坦右 \(\mathbf{A}_{\alpha}\) -模。证明：若诸 \(\mathbf{A}_{\alpha}\) 都右凝聚，则 \(\mathbf{A}\) 亦然。（注意 \(\mathbf{A}\) 是平坦 \(\mathbf{A}_{\alpha}\) -模（对所有 \(\alpha\) ），且若 \(\mathbf{E}\) 是 \(\mathbf{A}_d\) 的有限生成子模，则存在指标 \(\alpha\) 以及有限生成子模 \(\mathbf{E}_{\alpha}\) （属于 \((\mathbf{A}_{\alpha})_d\) ），使 \(\mathbf{E}_{\alpha} \otimes_{\mathbf{A}_{\alpha}} \mathbf{A}\) 同构于 \(\mathbf{E}\) 。）
\end{enumerate}

* (f) 由 (e) 推出：Noether 交换环上的每个多项式环（在任何有限或无限未定元集上）都是凝聚的。由此推出：凝聚环的商环未必凝聚。*

\begin{enumerate}
\def\labelenumi{(\alph{enumi})}
\setcounter{enumi}{6}
\tightlist
\item
  为使 A 左凝聚，必须且只需 A 的每个元素的左零化子是有限生成的，且 A 的两个有限生成左理想的交是有限生成的（利用习题 6(f)）。
\end{enumerate}

¶ 13. 设 A、B 是两个环，F 是 (A, B)-双模，G 是右 B-模。证明：若 G 是内射的（《代数学》第二章 §2 习题 11）且 F 是平坦左 A-模，则右 A-模 \(\operatorname{Hom}_{B}(F, G)\) 是内射的。（利用同构

\[
\operatorname{Hom} _ {\mathbf {A}} (\mathbf {E}, \operatorname{Hom} _ {\mathbf {B}} (\mathbf {F}, \mathbf {G})) \rightarrow \operatorname{Hom} _ {\mathbf {B}} (\mathbf {E} \otimes_ {\mathbf {A}} \mathbf {F}, \mathbf {G})
\]

其中 E 是右 A-模（《代数学》第二章 §4 第 1 目）。）

¶ 14. 设 A、B 是两个环，E 是左 A-模，F 是 (A, B)-双模，G 是右 B-模；考虑典范同态（《代数学》第二章 §4 习题 5）

\[
\sigma \colon \operatorname{Hom} _ {\mathbb {B}} (\mathrm{F}, \mathrm{G}) \otimes_ {\mathrm{A}} \mathrm{E} \rightarrow \operatorname{Hom} _ {\mathbb {B}} (\operatorname{Hom} _ {\mathrm{A}} (\mathrm{E}, \mathrm{F}), \mathrm{G})
\]

它满足 \((\sigma(u \otimes x))(v) = u(v(x))\) ，对所有 \(x \in \mathbf{E}\) ， \(u \in \operatorname{Hom}_{\mathbf{B}}(\mathbf{F}, \mathbf{G})\) ，

\[
v \in \operatorname{Hom} _ {\mathbf {A}} (\mathrm{E}, \mathbf {F}).
\]

证明：若 G 是内射 B-模（《代数学》第二章 §2 习题 11）且 E 有限表示，则 \(\sigma\) 是双射。（先考虑 E 是自由且有限生成的情形。）

¶ 15. 设 A 是环。证明每个平坦且有限表示的左 A-模 E 都是投射的。（给定左 A-模的满同态 \(u: \mathbf{F} \to \mathbf{F}'\) ，设 \(u'\) 是同态

\[
\operatorname{Hom} \left(1 _ {\mathbf {E}}, u\right): \operatorname{Hom} _ {\mathbf {A}} (\mathrm{E}, \mathrm{F}) \rightarrow \operatorname{Hom} _ {\mathbf {A}} (\mathrm{E}, \mathrm{F} ^ {\prime \prime})
\]

且设 \(\bar{u}\) 是同态

\[
\operatorname{Hom} \left(u ^ {\prime}, 1 _ {G}\right): \operatorname{Hom} _ {\mathbf {Z}} \left(\operatorname{Hom} _ {A} (E, F ^ {\prime \prime}), G\right)\rightarrow \operatorname{Hom} _ {\mathbf {Z}} \left(\operatorname{Hom} _ {A} (E, F), G\right),
\]

其中 G 是可除 Z-模。先利用习题 14 证明 \(\bar{u}\) 是单射；然后取适当的 G（《代数学》第二章 §2 习题 14），证明 \(u'\) 是满射。）

\begin{enumerate}
\def\labelenumi{\arabic{enumi}.}
\setcounter{enumi}{15}
\tightlist
\item
  设 A 是环， \(a\) 是 A 的元素。证明下列性质等价。
\end{enumerate}

\((\alpha)\) \(a \in aAa\) 。

(β) aA 是模 \(A_{d}\) 的直因子。

(γ) \(A_{d}/aA\) 是平坦右 A-模。

(δ) 对 A 的每个左理想 b， \(aA \cap b = ab\) 。

（利用第 6 目命题 7 的推论证明 \((\gamma)\) 与 \((\delta)\) 的等价性，并直接证明 \((\delta)\) 蕴含 \((\alpha)\) 且 \((\alpha)\) 蕴含 \((\beta)\) ，办法是证明存在幂等元 \(e \in a\mathbf{A}\) 使 \(e\mathbf{A} = a\mathbf{A}\) 。）

\begin{enumerate}
\def\labelenumi{\arabic{enumi}.}
\setcounter{enumi}{16}
\tightlist
\item
  设 A 是环。证明下列性质等价：
\end{enumerate}

(α) 每个元素 \(a \in A\) 都满足习题 16 的等价性质。

(β) A 的每个有限生成右理想都是 \(A_{d}\) 的直因子。

(γ) 每个左 A-模都是平坦的。

(δ) 每个右 A-模都是平坦的。

这时称 A 为绝对平坦环 (*)。

（为看出 \((\alpha)\) 蕴含 \((\beta)\) ，利用《代数学》第八章 §6 习题 15(b)。）

¶ 18. 设 A 是绝对平坦环（习题 17）。

\begin{enumerate}
\def\labelenumi{(\alph{enumi})}
\item
  设 P 是投射右 A-模。证明 P 的每个有限生成子模 E 都是 P 的直因子。（归结为 P 自由且有限生成的情形。然后注意 P/E 是有限表示的并利用习题 15。）
\item
  证明每个投射右 A-模 P 都是单生成子模的直和，这些子模同构于 A 的单生成右理想。（利用 Kaplansky 定理（《代数学》第二章 §2 习题 3）把问题归结为 P 由可数族元素生成的情形，然后利用 (a)。）
\item
给出绝对平坦环 A 与非投射的有限生成 A-模的一个例子（考虑 A 对一个非有限生成理想的商；参看《代数学》第八章 §6 习题 15(f) 与《交换代数》第二章 §4 习题 17）。
\end{enumerate}

\begin{enumerate}
\def\labelenumi{\arabic{enumi}.}
\setcounter{enumi}{18}
\tightlist
\item
  设 A 是环。证明下列性质等价：
\end{enumerate}

\((\alpha)\) A 是半单的。

(β) A 的每个右理想都是内射 A-模。

(γ) 每个右 A-模都是投射的。

(δ) 每个右 A-模都是内射的。

\begin{enumerate}
\def\labelenumi{\arabic{enumi}.}
\setcounter{enumi}{19}
\item
  设 A 是整环，B 是作为 A-模平坦的 A-代数，M 是无挠 A-模。证明：若 \(t \in B\) 不是零因子，则对 \(z \in B \otimes_{A} M\) 关系 t. z = 0 蕴含 z = 0。（归结为 M 有限生成的情形，且通过把 M 嵌入一个有限生成的自由 A-模，归结为 M = A 的情形。）
\item
  设 S 是交换环，R 是交换 S-代数，B 是 S-代数（交换或非交换）， \(\mathrm{B}_{(\mathrm{R})}\) 是由 B 经纯量扩张得到的 R-代数。假设 R 是平坦 S-模且 B 是有限生成 S-模。若 Z 是 B 的中心，证明从 \(Z_{(R)} = Z \otimes_{S} R\) 到 \(\mathrm{B}_{(\mathrm{R})}\) 的典范同态是 \(Z_{(R)}\) 到 \(\mathrm{B}_{(\mathrm{R})}\) 的中心上的同构。（利用正合序列
\end{enumerate}

\[
0 \longrightarrow \mathrm{Z} \longrightarrow \mathrm{B} \stackrel {\theta} {\longrightarrow} \operatorname{Hom} _ {\mathrm{s}} (\mathrm{B}, \mathrm{B})
\]

其中 \(\theta (x)(y) = xy - yx\) ，以及第 10 目的命题 11。）

\begin{enumerate}
\def\labelenumi{\arabic{enumi}.}
\setcounter{enumi}{21}
\tightlist
\item
  设 E 是左 A-模。对 A 的每个右理想 a 与每个元素 \(a \in A\) ，用 a: a 表示 \(x \in A\) （满足 \(ax \in a\) ）所成的集合，用 aE: a 表示 \(y \in E\) （满足 \(ay \in aE\) ）所成的集合。这时显然有 (a: a)E ⊂ aE: a。证明：为使 E 平坦，必须且只需对 A 的每个右理想 a 与每个元素 \(a \in A\) ，(a: a)E = aE: a。（为看出条件是必要的，考虑右 A-模的正合序列 \(0 \to (a: a)/a \stackrel{\psi}{\to} A_d/a \stackrel{\varphi}{\to} A_d/a\) ，其中 \(\psi\) 是典范单射， \(\phi\) 是由用 a 左乘再取商得到的映射。为看出条件是充分的，应用第 11 目命题 13 的推论 1 的判别准则：从一个关系 \(\sum_{i=1}^{n} a_i x_i = 0\) 出发，其中 \(a_i \in A\) ， \(x_i \in E\) ，把假设应用于理想
\end{enumerate}

\(a_{2}=\sum_{i=2}^{n}a_{i}A\) 与元素 \(a_{1}\) ，并对 n 用归纳法论证。）

\begin{enumerate}
\def\labelenumi{\arabic{enumi}.}
\setcounter{enumi}{22}
\tightlist
\item
  \begin{enumerate}
  \def\labelenumii{(\alph{enumii})}
  \tightlist
  \item
  设 \(0 \to \mathbf{R} \to \mathbf{L} \to \mathbf{E} \to 0\) 是左 A-模的正合序列，其中 \(\mathbf{L}\) 是自由 A-模；设 \((e_{\alpha})\) 是 \(\mathbf{L}\) 的基。证明下列条件等价：
  \end{enumerate}
\end{enumerate}

(α) E 是平坦的。

(β) 对所有 \(x \in \mathbb{R}\) ，若 \(\mathfrak{a}_x\) 是 \(x\) 关于基 \((e_{\alpha})\) 的分量生成的右理想，则 \(x \in \mathbb{R}\mathfrak{a}_x\) 。

\((\gamma)\) 对所有 \(x \in \mathbb{R}\) ，存在同态 \(u_x: \mathbf{L} \to \mathbf{R}\) 使 \(u_x(x) = x\) 。

\begin{enumerate}
\def\labelenumi{(\arabic{enumi})}
\setcounter{enumi}{7}
\tightlist
\item
对 R 中元素的每个有限序列 \((x_{i})_{1\leqslant i\leqslant n}\) ，存在同态 \(u\colon \mathbf{L}\to \mathbb{R}\) 使 \(u(x_{i}) = x_{i}\) （对 \(1\leqslant i\leqslant n\) ）。（利用第 6 目命题 7 的推论。）
\end{enumerate}

\begin{enumerate}
\def\labelenumi{(\alph{enumi})}
\setcounter{enumi}{1}
\tightlist
\item
  设 \(\mathfrak{a}\) 是 \(\mathbf{A}\) 的左理想，使 \(\mathbf{A} / \mathfrak{a}\) 是平坦 \(\mathbf{A}\) -模。证明：对每个有限生成左理想 \(\mathfrak{b} \subset \mathfrak{a}\) ，存在 \(x \in \mathbf{A}\) 使
\end{enumerate}

\[
\mathfrak {b} \subset \mathrm{A} x \subset \mathfrak {a}
\]

（利用 (a) 的条件 (δ)）。

\begin{enumerate}
\def\labelenumi{(\alph{enumi})}
\setcounter{enumi}{2}
\item
  由 (a) 导出习题 15 的结果的一个新证明。
\item
  设 \(\mathfrak{r}\) 是 \(\mathbf{A}\) 的 Jacobson 根，且 \(0 \to \mathbf{R} \to \mathbf{L} \to \mathbf{E} \to 0\) 是左 A-模的正合序列，其中 \(\mathbf{L}\) 是自由的。假设 \(\mathbf{E}\) 平坦且 \(\mathbf{R}\) 包含于 \(\mathfrak{rL}\) 。证明 \(\mathbf{R} = 0\) （用 (a) 的记号，注意 \(\mathfrak{a}_x\) 是有限生成理想且 \(\mathfrak{a}_x = \mathfrak{a}_x\mathfrak{r}\) ）。
\item
  设 E 是有限生成的平坦 A-模；假设存在 A 的包含于 Jacobson 根中的双边理想 b，使 E/bE 是自由 (A/b)-模。证明这时 E 是自由 A-模（注意存在有限生成自由 A-模 L 使 L/bL 同构于 E/bE，并利用《代数学》第八章 §6 第 3 目命题 6 的推论 4；然后应用 (d)）。
\end{enumerate}

\begin{enumerate}
\def\labelenumi{\arabic{enumi}.}
\setcounter{enumi}{23}
\tightlist
\item
右 A-模 M 的子模 \(M'\) 称为纯的，如果用 \(j: M' \rightarrow M\) 表示典范单射，同态
\end{enumerate}

\[
j \otimes 1 _ {N} \colon M ^ {\prime} \otimes_ {A} N \rightarrow M \otimes_ {A} N
\]

对每个左 A-模 N 是单射。若 \(M'\) 是 M 的直因子或 \(M/M'\) 平坦，情形就是如此，但这两个条件不是必要的（习题 4）。

\begin{enumerate}
\def\labelenumi{(\alph{enumi})}
\item
  证明：为使 \(\mathbf{M}'\) 是 \(\mathbf{M}\) 的纯子模，必须且只需：若 \((m_i')_{i\in I}\) 是 \(\mathbf{M}'\) 的元素的有限族， \((x_j)_{j\in J}\) 是 \(\mathbf{M}\) 的元素的族，使 \(m_i' = \sum_{j\in J} x_j a_{jl}\) 对所有 \(i\in I\) 成立，且 \((a_{jl})\) 是 \(\mathbf{A}\) 的元素的族，则存在族 \((x_j')_{j\in J}\) （元素属于 \(\mathbf{M}'\) ）使 \(m_i' = \sum_{j\in J} x_j' a_{jl}\) 对所有 \(i\in I\) 成立。（为看出条件是充分的，利用第 11 目的引理 10 证明 \(\mathbf{M}' \otimes_{\mathbf{A}} \mathbf{N} \to \mathbf{M} \otimes_{\mathbf{A}} \mathbf{N}\) 对每个有限生成左 A-模 \(\mathbf{N}\) 是单射；为看出条件是必要的，考虑有限生成左 A-模 \(\mathbf{N} = \mathbf{L}/\mathbf{R}\) ，其中 \(\mathbf{L}\) 是有限生成自由 A-模而 \(\mathbf{R}\) 是 \(\mathbf{L}\) 的有限生成子模。）由这个判别准则推出：若 \(\mathbf{A}\) 是主理想整环，则 A-模的纯子模概念与《代数学》第七章 §2 习题 7 中的概念一致。
\item
  设 M 是右 A-模，M' 是 M 的子模，M'' 是 M' 的子模。证明：若 M' 是 M 的纯子模且 M'' 是 M' 的纯子模，则 M'' 是 M 的纯子模且 M'/M'' 是 M/M'' 的纯子模。若 M'' 是 M 的纯子模，则 M'' 是 M' 的纯子模。
\item
  证明：若 N 与 P 是 M 的两个子模，使 \(N \cap P\) 与 \(N + P\) 在 M 中是纯的，则 N 与 P 是 M 的纯子模。给出 \(Z^{2}\) 的两个子模 N、P 的例子，它们在 \(Z^{2}\) 中是纯的，但 \(N + P\) 不是 \(Z^{2}\) 的纯子模。
\item
  设 C 是交换环，E、F 是两个 C-模；证明：若 \(E'\) （相应地 \(F'\) ）是 E（相应地 F）的纯子模，则典范映射 \(E' \otimes_{C} F' \to E \otimes_{C} F\) 是单射，且把 \(E' \otimes_{C} F'\) 等同于 \(E \otimes_{C} F\) 的一个纯子模。
\item
  设 \(\rho: \mathbf{A} \to \mathbf{B}\) 是环同态， \(\mathbf{M}\) 是右 A-模， \(\mathbf{M}'\) 是 \(\mathbf{M}\) 的纯子模。证明 \(\mathbf{M}_{(\mathbf{B})}' = \mathbf{M}' \otimes_{\mathbf{A}} \mathbf{B}\) 典范地等同于 \(\mathbf{M}_{(\mathbf{B})} = \mathbf{M} \otimes_{\mathbf{A}} \mathbf{B}\) 的一个纯子模。
\end{enumerate}

